\begin{abstract}
Robot manipulation requires experience that explains which interaction phase is active, what an action achieved, and which decision its feedback supports. We propose RoboRAG, a framework that carries this information from recorded episodes into plans for unseen task compositions. Its central hypothesis is that preserving phase conditions and observed outcomes throughout this process improves experience selection and plan following. An annotation pipeline constructs temporal interaction graphs with phase-instance nodes and typed edges for prerequisites, conditional transitions, and recovery. A versioned bank supports retrieval of complete procedures and local subgraphs. A GPT planner creates an initial plan under the shared schema, then uses retrieved experience to revise its object bindings, parameters, and structure. A lower-cost supervisor follows approved phases and recovery branches; a gateway and local controllers check and execute its requests. The evaluation compares graph and sequence methods with matched evidence, tests adaptation and delegated execution, and separates interface reuse from MCP transport across proposed simulation and physical backends. This design-stage manuscript specifies the method and evaluation protocol; available pilots concern an earlier chunk-based memory controller.
\end{abstract}

% ABSTRACT_SUMMARY: RoboRAG proposes a shared phase-graph schema connecting episode annotation, versioned graph-bank retrieval, task-conditioned plan adaptation, and delegated execution. Matched evidence, representation, adaptation, execution, and interface comparisons test the specific hypothesis that preserving phase conditions and observed outcomes improves experience use. This is a design-stage manuscript; existing pilots concern a predecessor chunk-memory controller.
